\section{Concrete Example: Length-Map Fusion}
\label{sec:example}

To illustrate the correspondence between supercompilation and cyclic proof search, we present a worked example: the classic deforestation problem of fusing \texttt{length} and \texttt{map}.

\subsection{The input term}

Consider the composed function that applies a function to every element of a list, then computes the length:
\[
\lambda f.\, \lambda l.\, \mathtt{length}\,(\mathtt{map}\,f\,l)
\]

This creates an intermediate list that is immediately consumed, which is inefficient. The expected optimized output is simply:
\[
\lambda f.\, \lambda l.\, \mathtt{length}\,l
\]

\subsection{Supercompilation trace}

The supercompiler proceeds as follows:

\begin{enumerate}
\item \textbf{Initial configuration:} $\Gamma \vdash \lambda f.\, \lambda l.\, \mathtt{length}\,(\mathtt{map}\,f\,l) : (\mathtt{Nat} \to \mathtt{Nat}) \to \mathtt{List} \to \mathtt{Nat}$

\item \textbf{Drive (unfold):} Unfold \texttt{length} and \texttt{map} fixpoint definitions, exposing:
\[
\lambda f.\, \lambda l.\, \langle l \mid \mathtt{nil} \to \cdots;\; \mathtt{cons}\,x\,xs \to \cdots \rangle
\]

\item \textbf{Split (case analysis):} Split on the neutral variable $l$, creating two branches:
\begin{itemize}
  \item \textbf{Nil branch:} $\mathtt{length}\,(\mathtt{map}\,f\,\mathtt{nil}) \leadsto 0$
  \item \textbf{Cons branch:} $\mathtt{length}\,(\mathtt{map}\,f\,(\mathtt{cons}\,x\,xs))$
\end{itemize}

\item \textbf{Drive (beta-iota):} In the cons branch, drive performs:
\[
\mathtt{length}\,(\mathtt{cons}\,(f\,x)\,(\mathtt{map}\,f\,xs)) \leadsto \mathtt{succ}\,(\mathtt{length}\,(\mathtt{map}\,f\,xs))
\]

\item \textbf{Fold (memoization):} Recognize that $\mathtt{length}\,(\mathtt{map}\,f\,xs)$ is a recursive call matching the initial configuration. Create a backlink to the root.

\item \textbf{Result:} The residual term is:
\[
\lambda f.\, \lambda l.\, \langle l \mid \mathtt{nil} \to 0;\; \mathtt{cons}\,x\,xs \to \mathtt{succ}\,(\mathtt{REC}\,f\,xs) \rangle
\]
where $\mathtt{REC}$ is the recursive call.
\end{enumerate}

Note that the intermediate list construction \texttt{map f l} never appears in the residual term---it has been fused away.

\subsection{Corresponding cyclic proof}

The supercompilation trace above corresponds to a cyclic sequent proof with the following structure:

\begin{enumerate}
\item \textbf{Root node:} Sequent $\Gamma \vdash \lambda f.\, \lambda l.\, \mathtt{length}\,(\mathtt{map}\,f\,l) : (\mathtt{Nat} \to \mathtt{Nat}) \to \mathtt{List} \to \mathtt{Nat}$

\item \textbf{Asynchronous edges:} From unfolding \texttt{length} and \texttt{map}, each drive step creates a single successor with the driven term.

\item \textbf{Synchronous branching:} The case split on $l$ creates two premises (nil and cons branches).

\item \textbf{Cyclic backlink:} The fold creates an edge back to the root node, forming a cycle.

\item \textbf{Local validity:} Each node satisfies its sequent rule:
\begin{itemize}
  \item Driving steps satisfy \texttt{dr\_cbn\_once}
  \item The case split satisfies \texttt{dr\_split\_case\_var}
  \item The backlink is justified by configuration equality
\end{itemize}
\end{enumerate}

\subsection{Graph structure}

The resulting proof graph has three nodes: the root (initial configuration), a nil branch returning 0, and a cons branch performing a recursive call. The cons branch has a backlink to the root, forming the cycle.

The cycle (from the cons branch back to root) represents the recursive structure. The \emph{trace condition} for this cycle is satisfied because the case-split is a \emph{progress point}: the recursive call is on $xs$, which is structurally smaller than $l$.

This example is mechanised in Rocq; the supercompiler successfully processes the input and the step correspondence theorems ensure a matching cyclic sequent proof. Additional benchmarks (append-length, append-associativity via 2-pass supercompilation, take-drop fusion, and vector-index normalisation) are validated by computational reflection.
